\section{Uniform contraction of contiguous blocks}
\label{sec:ldp_inhomogeneous_contraction}

A quantitative condition on the transfer maps of actual contiguous blocks
provides a uniform positive-part rate. The condition below is Choi domination,
which gives the quantum Doeblin contraction of~\cite[Theorem~8.17]{Wolf2012Quantum}.
It is an additional structural hypothesis; it is not inferred from the
inhomogeneous finite-correlation definition of~\cite{Malz2024Preparation}.
In particular, separate spectral gaps of the microscopic site maps do not
control ordered products of different maps.

\begin{theorem}[Transfer map of an inhomogeneous block]
    \label{thm:ldp_chain_transfer_product}
    \lean{MPSChainTensor.prod_transferMap_apply,
        MPSChainTensor.transferMap_blockTensor,
        MPSChainTensor.trace_transferMap_blockTensor}
    \leanok
    \uses{def:ldp_chain_block_tensor}
    For a chain $A_0,\dots,A_{n-1}$ with common bond dimension $D$, let $B$
    be its blocked tensor. Then
    \begin{align}
        E_B(X)
         & =\sum_{i_0,\dots,i_{n-1}}
        A_0^{i_0}\cdots A_{n-1}^{i_{n-1}}X
        (A_0^{i_0}\cdots A_{n-1}^{i_{n-1}})^\dagger
        \notag                                        \\
         & =(E_{A_0}\circ\cdots\circ E_{A_{n-1}})(X).
    \end{align}
    The first site's map is outermost. If every site map preserves trace,
    so does $E_B$.
\end{theorem}
\begin{proof}\leanok
    Expand the Kraus sum, split off the first physical index, and proceed
    by induction on $n$. The two adjacent matrix products have the same
    order as the composition; the adjoints appear in the reverse order.
    Trace preservation follows by applying it to the outermost map at
    each step.
\end{proof}

\begin{theorem}[Uniform Doeblin contraction along a chain]
    \label{thm:ldp_chain_doeblin}
    \lean{MPSChainTensor.traceNorm_transferMap_blockTensor_sub_le_of_choi_domination}
    \leanok
    \uses{thm:ldp_chain_transfer_product}
    Let $\sigma$ be a density matrix and $0\le\eta\le1$. Suppose every
    site's transfer map preserves trace and its normalized Choi matrix
    satisfies $J(E_{A_i})\ge(\eta/D)(\sigma\otimes I)$. For the blocked
    tensor $B$ and any two density matrices $\rho_1,\rho_2$,
    \begin{align}
        \lVert E_B(\rho_1)-E_B(\rho_2)\rVert_1
         & \le(1-\eta)^n\lVert\rho_1-\rho_2\rVert_1.
    \end{align}
    The sites may themselves be contiguous blocks, provided the hypothesis
    is imposed on their actual transfer maps.
\end{theorem}
\begin{proof}\leanok
    Apply the Choi-dominated quantum Doeblin theorem to the first map and
    the two density matrices produced by the remaining maps. Induction
    gives one factor $1-\eta$ for each actual map in the ordered product.
\end{proof}

\begin{theorem}[Uniform Gram and positive-part bounds]
    \label{thm:ldp_chain_positive_part_rate}
    \lean{MPSPreparation.gram_eq_smul_choi_transpose,
        MPSPreparation.exists_norm_polarPos_chainBlockTensor_sub_le_of_choi_domination}
    \leanok
    \uses{thm:ldp_chain_transfer_product, def:ldp_chain_block_tensor}
    Fix $\sigma>0$ with $\operatorname{Tr}\sigma=1$. There is $K$, depending
    only on $D$ and $\sigma$, such that the following holds. Partition a chain
    of $N$ sites into contiguous blocks of lengths $\ell_j\le s$, with
    $s\ge1$, and let $F_j$ be their transfer maps. Suppose each $F_j$ preserves
    trace, fixes $\sigma$, and satisfies
    $J(F_j)\ge(\eta/D)(\sigma\otimes I)$ for one $0\le\eta\le1$.
    For the physical matrix $B$ of the whole chain and its positive polar
    factor $P_B$, the operator-norm bounds are
    \begin{align}
        \lVert B^\dagger B-\sigma^\mathsf{T}\otimes I\rVert
         & \le 2D(1-\eta)^{\lfloor N/s\rfloor},\notag \\
        \lVert P_B-(\sqrt\sigma)^\mathsf{T}\otimes I\rVert
         & \le K(1-\eta)^{\lfloor N/s\rfloor}.
    \end{align}
    The constants are independent of the site tensors, the physical
    dimension, the chain length, and the partition. For fixed $s$ and
    $0<\eta<1$, these are uniform exponential bounds in physical length.
\end{theorem}
\begin{proof}\leanok
    Put $P_\sigma(X)=\operatorname{Tr}(X)\sigma$ and
    $R_j=F_j-\eta P_\sigma$. Choi domination makes every $R_j$ completely
    positive. Its trace is scaled by $1-\eta$. The common fixed state and
    trace preservation give, for $m$ contiguous blocks,
    \begin{align}
        F_0\circ\cdots\circ F_{m-1}
         & =[1-(1-\eta)^m]P_\sigma+R_0\circ\cdots\circ R_{m-1}.
    \end{align}
    The normalized Choi matrix of the residual product is positive and
    has trace $(1-\eta)^m$. A positive matrix has operator norm at most its
    trace. Also $J(P_\sigma)$ is a density matrix. Since
    $B^\dagger B=D J(E_B)^\mathsf{T}$, the triangle inequality gives the Gram
    bound $2D(1-\eta)^m$. The partition implies $N\le ms$, hence
    $m\ge\lfloor N/s\rfloor$.

    The limiting Gram matrix $\sigma^\mathsf{T}\otimes I$ is positive
    definite. Apply the square-root Lipschitz estimate at that fixed
    matrix, with constant $L_\sigma$, and take $K=2DL_\sigma$.
\end{proof}

These results do not by themselves prove the uniform pair-approximation
hypothesis for general inhomogeneous chains. A common faithful fixed state
and quantitative domination on actual contiguous blocks are assumed here.
Varying rectangular bond spaces and varying fixed states require additional
arguments. Section~\ref{sec:ldp_inhomogeneous_mixing} proves the overlap and
normalization bounds under a uniform ordered-transfer mixing hypothesis;
qualitative finite correlation alone still does not supply that hypothesis
\cite{gap:mswc24_inhomogeneous_scope}.
